\documentclass[11pt]{article}

\usepackage[utf8]{inputenc}
\usepackage[T1]{fontenc}
\usepackage{lmodern}
\usepackage{microtype}
\usepackage[margin=1in]{geometry}
\usepackage{amsmath,amssymb}
\usepackage{booktabs}
\usepackage{enumitem}
\usepackage{xcolor}
\usepackage[hidelinks]{hyperref}
\usepackage{siunitx}

\sisetup{detect-weight=true}

\title{\textbf{Building the AI We Actually Need}\\[6pt]
\large A Global Photonic Network for Cooperative, Problem-Focused Development}

\author{%
  David O'Brien\thanks{Corresponding author. Affiliated with the Georgia Tech
  AI Safety Initiative (AISI). A working pre-print circulated for feedback and
  collaboration. Drafting and synthesis assistance from a large language model
  (Anthropic Claude); all claims and errors are the author's responsibility.}\\
  \small Georgia Tech AI Safety Initiative (AISI)
}

\date{Version 0.4 --- \today}

\begin{document}
\maketitle

\begin{abstract}
\noindent Powerful AI will almost certainly be built. The question is whether it
will emerge from a competitive race---with resources wasted on redundancy,
secrecy, and mutual defensive positioning---or from a cooperative effort pooling
human and computational resources toward actually \emph{solving problems}.
\par
Until now, the bottleneck preventing planetary-scale cooperation was physical:
moving neural network activations between distant compute sites consumed more
energy and latency than made global pooling practical. That bottleneck is now
removable. This paper proposes a global photonic transport fabric---built on
existing fiber infrastructure---that solves the data-movement problem. With it
solved, a single, globally distributed AI system becomes not a distant hope but
an engineering project: the fastest way to frontier capability, the least
wasteful deployment of human and computational effort, and the configuration
most able to actually help with the problems that matter---climate, disease,
energy, scientific discovery, human flourishing.
\par
We show that the technical barrier has no unsolved problems: every component
exists as a commercial product or mature prototype. We give a proposal-depth
architecture, performance estimates, a realistic three-phase roadmap (Phase~1 in
2--3 years, \$35--50M), and the concrete mechanisms that would keep cooperation
stable and prevent capture. We address this to builders, researchers, and the
institutions that fund them: not as a moral appeal, but as the engineering path
that is simultaneously the fastest, the least wasteful, and the most useful.

\vspace{8pt}
\noindent\textbf{Keywords:} cooperative AI development, distributed compute,
photonic transport, data-center interconnects, global infrastructure, problem-solving
capability, resource efficiency.
\end{abstract}

\clearpage
\tableofcontents
\clearpage

% ======================================================================
\section{Why We Should Want to Build This Together}
\label{sec:why}

\subsection{The race is expensive and slow}

Today, frontier AI development is driven by competition. Multiple labs, multiple
firms, multiple nations are racing to build capable systems. This is
extraordinarily wasteful:

\begin{itemize}[leftmargin=1.6em,itemsep=4pt]
  \item \textbf{Redundant effort.} Teams solve the same problems in parallel,
  independently, without learning from each other. The compute spent on solving
  training techniques twice, three times, four times over is compute not spent
  on making systems safer, more capable, or more useful.
  \item \textbf{Secrecy overhead.} Each lab hoards insights, methods, even
  infrastructure designs. The tax of not sharing is enormous: algorithms,
  hardware designs, and datasets that could be public knowledge are kept private.
  \item \textbf{Defensive positioning.} Resources are diverted into securing
  talent, locking down supply chains, and positioning for advantage against
  rivals. In the worst case, safety work is deferred because a lab fears being
  overtaken. In the best case, it's just inefficient.
  \item \textbf{Fragmentation of compute.} No single actor can pool all
  available compute. Each builds their own stack in isolation. The result is
  that we have many smaller systems rather than one larger one---and smaller
  systems are slower to solve hard problems.
\end{itemize}

\noindent By contrast, a cooperative effort would:
\begin{itemize}[leftmargin=1.6em,itemsep=4pt]
  \item Pool all available compute and talent into one focused effort.
  \item Eliminate redundancy: solve problems once, share solutions immediately.
  \item Enable transparency and openness at every stage, speeding iteration and
  learning.
  \item Concentrate resources on what matters: capability and useful application,
  not competitive positioning.
\end{itemize}

\subsection{The most powerful system is a cooperative one}

There is a simple arithmetic here. Frontier capability scales with compute. The
most powerful system we can build is therefore the one that aggregates the most
compute. No single firm, lab, or nation controls all the compute that matters.
A coalition of all of them does.

\textbf{Hence: the fastest path to the most capable AI is a cooperative global
effort.} Not because cooperation is virtuous, but because it is the only way to
aggregate enough resources to win.

\subsection{Cooperation enables the safety budget we actually need}

There is a secondary but crucial point. A system built in a race must ship when
it is competitive---that is, as soon as it is better than rivals. There is no
margin for extra time spent on alignment, on testing, on building in
interpretability and safeguards. Speed is the primary selection pressure.

A system built cooperatively does not face that pressure. If all major actors are
in the coalition, there is no ``being overtaken.'' There is time to:
\begin{itemize}[leftmargin=1.6em,itemsep=4pt]
  \item Build alignment and interpretability in from the start, rather than
  bolting it on.
  \item Run extensive testing and red-teaming without fear of a deadline.
  \item Implement governance, transparency, and benefit-sharing mechanisms.
  \item Move deliberately rather than at the speed of fear.
\end{itemize}

A cooperative system can be built \emph{right} because there is no race penalty
for taking the time.

\subsection{The bottleneck: we know what it is, and we can solve it}

Until now, the reason planetary-scale cooperation has seemed impractical is
physical. Moving neural network activations between distant compute sites is
slow (hundreds of milliseconds of fiber latency plus electronic switching delay)
and energy-intensive (the data-movement overhead dominates the arithmetic). This
makes a single global system infeasible.

That bottleneck is now solvable. The solution is in Section~\ref{sec:bottleneck}.
With it solved, cooperative global compute pooling becomes not just possible but
pragmatic.

% ======================================================================
\section{The Bottleneck: Data Movement, and How to Remove It}
\label{sec:bottleneck}

\subsection{Why data movement is the constraint}

In training or inference on large models, the limiting factor is not
computation---we have plenty of compute, and it is cheap. The limiting factor is
moving activations between compute nodes. For a distributed system spanning
continents:
\begin{itemize}[leftmargin=1.6em,itemsep=4pt]
  \item \textbf{Latency.} Electronic packet-switched networks impose 150~ms of
  fiber propagation plus 200--250~ms of switching/queueing overhead. For
  distributed inference (where model layers are partitioned across sites), this
  becomes a hard constraint.
  \item \textbf{Energy.} Moving a bit across a data center via electronic
  switching costs roughly \SI{1}{pJ}. For a petabyte-scale inference fleet,
  this dominates the total energy budget. The compute itself is almost free by
  comparison.
\end{itemize}

\noindent These constraints prevent the kind of planetary-scale compute pooling that would
make cooperation obviously efficient.

\subsection{The core idea: optical transport}

Fiber-optic cables already span the planet---roughly \SI{10e6}{km} of submarine
fiber, much of it dark (underutilized). Dense wavelength-division multiplexing
(WDM) routinely carries 50--400 independent signals per fiber, each at 10--100~Gbps.

The idea is simple: instead of converting activations to electronic packets,
routing them electronically, and converting back to optical, encode them
\emph{directly as optical signals} and route them by wavelength. This yields:

\begin{itemize}[leftmargin=1.6em,itemsep=4pt]
  \item \textbf{Energy reduction:} Optical transport costs roughly \SI{0.01}{pJ/bit}.
  A \(100\times\) improvement in the dominant energy cost of global compute pooling.
  \item \textbf{Latency reduction:} Eliminate electronic switching. Fiber
  propagation is immutable (\SI{30}{ms} intercontinental), but switching overhead
  vanishes. Single-hop latency drops from \(\sim400\) ms to \(\sim30\) ms.
  \item \textbf{Practical feasibility:} Passive optical routing requires no
  per-signal active electronics. Wavelength demultiplexers are simple, proven,
  and cheap. This is not a laboratory concept; it is deployed infrastructure.
\end{itemize}

This is not optical \emph{computing} (where photons do the arithmetic).
Computation stays electronic at each node. Only \emph{transport} is optical---a
clean, modular separation that avoids the precision and scalability challenges
of analog photonic compute while capturing the energy and latency wins of
eliminating electronic switching.

% ======================================================================
\section{What It Would Look Like: Technical Architecture}
\label{sec:tech}

\subsection{System overview}

At each regional hub:
\begin{enumerate}[leftmargin=1.6em,itemsep=2pt]
  \item Inference produces 16-bit activations (electronic).
  \item Electro-optic modulators encode activations as intensity-modulated
  \SI{1550}{nm} light.
  \item A wavelength multiplexer combines 50--100 wavelengths onto fiber.
  \item Signals propagate through existing submarine and terrestrial cables.
  \item Erbium-doped fiber amplifiers (EDFAs) maintain signal power every
  \(\sim\!80\) km.
  \item At the destination, a demultiplexer, photodiodes, and 16-bit ADCs
  recover activations for the next layer.
  \item Electronic inference continues locally.
\end{enumerate}

\noindent The entire system ``lives'' in the electronic domain at each node. Optics is
purely transport.

\subsection{Signal representation and precision}

Activations are quantized to 16-bit unsigned integers. Modern neural networks
tolerate 8-bit (and often 4-bit) inference with $<1\%$ accuracy loss. Analysis of
the signal path shows effective precision of \(\sim\!13\)--14 bits after
quantization and accumulated noise---a comfortable margin.

\subsection{Long-haul transport and feasibility}

For a representative trans-Atlantic path (\(\sim\!6000\) km):
\begin{itemize}[leftmargin=1.6em,itemsep=2pt]
  \item Fiber attenuation: \(\sim\!0.2\) dB/km (standard telecom fiber).
  \item Amplifier spacing: EDFAs every \(\sim\!80\) km (\(\sim\!75\) spans total).
  \item Noise figure: \(\sim\!4\)--5 dB per amplifier; accumulated noise floor
  \(\sim\!-40\) dBm.
  \item Signal budget: transmit power \SI{+0}{dBm}; losses and gains balance
  to receiver sensitivity \SI{-20}{dBm}, with \SI{+6}{dB} margin for aging and
  drift.
  \item Propagation delay: \SI{30}{ms} one-way (speed of light in silica,
  irreducible).
  \item Channel isolation: WDM channels spaced at \SI{100}{GHz} (\(\sim\!0.4\) nm);
  adjacent-channel crosstalk $<-30$ dB (standard in telecom).
\end{itemize}

\noindent \textbf{Conclusion}: The signal budget is achievable with commercial-grade
components. No new physics required. No unsolved engineering problems.

\subsection{Wavelength routing and control}

Each datacenter pair is assigned a dedicated set of wavelengths. Activations for
a specific destination are encoded onto those wavelengths; at intermediate hubs,
passive optical add-drop multiplexers (OADMs) drop local traffic and pass through
to the next hop. A software-defined networking (SDN) controller manages
wavelength allocation and failover. Wavelengths are reused across geographically
disjoint paths (reuse factor \(\sim\!2\)--3$\times$).

A 10-hub global network needs \(\sim\!20\)--30 fiber pairs and 200--300 distinct
wavelengths. Modern WDM systems operate 400$+$ wavelengths; we use \(\sim\!50\)--75\% of
available capacity, leaving headroom for growth.

\subsection{Clock synchronization}

All nodes synchronize to a global reference (GPS-disciplined oscillator at major
hubs, refined by IEEE 1588v2 Precision Time Protocol over the optical network
itself). Achieves \(\pm\!10\) ns relative jitter---routine in telecom
infrastructure, requiring no new technology.

\subsection{Latency breakdown}

For a single intercontinental hop (US East Coast $\leftrightarrow$ Europe):
\begin{itemize}[leftmargin=1.6em,itemsep=2pt]
  \item Inference computation: \(\sim\!1\)--10 \(\mu\)s.
  \item Encoding (EOM + launch): \(\sim\!1\)--2 \(\mu\)s.
  \item Fiber propagation: \(\sim\!30\) ms (immutable).
  \item Reception (demux + photodetection + ADC): \(\sim\!1\)--2 \(\mu\)s.
  \item \textbf{Total}: \(\sim\!31\) ms.
\end{itemize}

\noindent For comparison, electronic inter-datacenter routing: \(\sim\!150\) ms fiber +
\(\sim\!200\)--250 ms switching/queueing = \(\sim\!350\)--400 ms. \textbf{A 10--13$\times$
improvement.}

For multi-hop distributed inference (three regions, three hops): \(\sim\!110\) ms
round-trip, making cross-continental model parallelism practical for batch and
many interactive workloads.

\subsection{Energy efficiency}

\begin{itemize}[leftmargin=1.6em,itemsep=2pt]
  \item Electronic inter-datacenter communication: \(\sim\!1\) pJ/bit.
  \item Optical transport: \(\sim\!0.01\) pJ/bit.
  \item For a large inference fleet (100M requests/day, 1M activations/request):
  reduces annual data-movement energy from \(\sim\!3.3\) GWh to \(\sim\!33\) MWh.
  \item As a fraction of total datacenter energy (compute + storage + cooling):
  saves \(\sim\!20\)--30\% under full deployment.
\end{itemize}

\subsection{Throughput}

Per wavelength: 16 Gbps (1 GHz clock $\times$ 16 bits). 50 wavelengths per fiber
pair: \(\sim\!800\) Gbps, or \(\sim\!50\)M activations/sec. A 10-pair global network
supports \(\sim\!500\)M activations/sec, sufficient for hundreds of thousands of
concurrent inferences.

% ======================================================================
% ======================================================================
\section{Technology Readiness: No Breakthroughs Needed}
\label{sec:trl}

Every component of this system already exists, either as a deployed commercial
product or as a mature research prototype:

\begin{description}[leftmargin=1.6em,itemsep=3pt]
  \item[Tier 1 \textemdash{} Off-the-shelf.] WDM infrastructure, 1550 nm fiber,
  EDFA amplifiers, Mach--Zehnder electro-optic modulators, photodiodes,
  transimpedance amplifiers, 16-bit ADCs, passive multiplexers/demultiplexers.
  Buyable today; suitable for a proof-of-concept testbed immediately.
  
  \item[Tier 2 \textemdash{} Emerging, needs optimization.] Cost and yield scaling of
  silicon photonic integrated circuits (10k$+$ channels already demonstrated in
  research labs). Monolithic integrated receivers (photodiode + TIA + ADC on one
  chip). Network-wide synchronization at \(\pm\!10\) ns. These are chip design
  and systems-integration problems, \emph{not} physics gaps. Development time:
  18--24 months with focused effort.
  
  \item[Tier 3 \textemdash{} Systems integration.] End-to-end protocol specification,
  wavelength-allocation control plane, fault detection and automatic failover,
  hardware/software co-design for latency guarantees. Non-trivial engineering,
  well within reach of a focused consortium with experience in distributed systems.
\end{description}

\noindent A Phase~1 trans-Atlantic testbed is estimated at 18--24 months of focused
engineering, with costs dominated by photonic-chip fabrication runs rather than
research. The barriers are organizational and economic, not scientific.

% ======================================================================
\section{How to Get There: A Phased Roadmap}
\label{sec:roadmap}

\subsection{Phase 1: Trans-Atlantic Proof of Concept (Years 1--3, \$35--50M)}

\textbf{Objective:} Demonstrate end-to-end optical activation transport at scale.

\textbf{Scope:}
\begin{itemize}[leftmargin=1.6em,itemsep=2pt]
  \item One leased dark-fiber pair on existing submarine cable (US East Coast
  $\leftrightarrow$ Europe).
  \item Two major datacenters (e.g., Virginia and Amsterdam).
  \item 10,000 input channels; 4,000 output channels (initial).
  \item Inference workload: distributed layers of a GPT-3 scale model
  (\(\sim\!175\)B parameters).
  \item Continuous operation: 6--12 months.
\end{itemize}

\textbf{Milestones:}
\begin{enumerate}[leftmargin=1.6em,itemsep=2pt]
  \item \textbf{Months 1--6: Subsystem design \& procurement.}
    \begin{itemize}[leftmargin=1.2em,itemsep=1pt]
      \item Procure and characterize lasers, modulators, detectors.
      \item Validate signal budget on lab testbed (\(\sim\!100\) m fiber).
      \item Finalize clock synchronization protocol.
    \end{itemize}
  
  \item \textbf{Months 6--12: System integration.}
    \begin{itemize}[leftmargin=1.2em,itemsep=1pt]
      \item Design and fabricate silicon photonic chip (10k channels).
      \item Integrate transceiver modules (lasers, EOMs, photodiodes, ADCs).
      \item Secure fiber connection to existing submarine cable landing station.
      \item Build SDN control plane for wavelength allocation.
    \end{itemize}
  
  \item \textbf{Months 12--18: Deployment \& characterization.}
    \begin{itemize}[leftmargin=1.2em,itemsep=1pt]
      \item Install hardware at both datacenters.
      \item Tune optical amplifiers, wavelength routing, synchronization.
      \item Run inference benchmarks: measure latency, throughput, bit-error rates.
      \item Validate energy reduction against electronic baseline.
    \end{itemize}
  
  \item \textbf{Months 18--24: Scaling \& optimization.}
    \begin{itemize}[leftmargin=1.2em,itemsep=1pt]
      \item Scale from 10k to 100k input channels (via additional wavelengths or higher
      modulation rates).
      \item Optimize for realistic inference patterns.
      \item Publish results; iterate design for Phase~2.
    \end{itemize}
\end{enumerate}

\textbf{Cost breakdown (\$35--50M):}
\begin{itemize}[leftmargin=1.6em,itemsep=2pt]
  \item Photonic chip design \& fabrication (2 runs): \$10M
  \item Transceiver hardware (lasers, EOMs, detectors, ADCs): \$5M
  \item Fiber lease \& amplifier infrastructure (5 years): \$8--10M
  \item Engineering labor (50 FTE over 3 years): \$12--15M
  \item Testing, characterization, contingency: \$3--5M
\end{itemize}

\textbf{Success criteria:}
\begin{itemize}[leftmargin=1.6em,itemsep=2pt]
  \item Sub-50 ms inter-datacenter activation transport latency
  \item $<1\%$ bit-error rate
  \item Demonstrated 50--100$\times$ energy reduction for interconnect vs.\ electronic baseline
  \item Sustained operation $>6$ months without significant degradation
  \item Integration with standard ML frameworks (PyTorch, JAX) demonstrated
  \item Published peer-reviewed results
\end{itemize}

\textbf{Risk level: Low.} Proven telecom technology; no new physics; existing components.

\subsection{Phase 2: Regional Hub Deployment (Years 4--6, \$150--200M)}

\textbf{Objective:} Extend to three regional hubs with full mesh connectivity and
production-grade reliability.

\textbf{Scope:}
\begin{itemize}[leftmargin=1.6em,itemsep=2pt]
  \item Three to four datacenters (US, Europe, Asia-Pacific) in partial mesh.
  \item Three fiber pairs (US $\leftrightarrow$ Europe, US $\leftrightarrow$ APAC,
  Europe $\leftrightarrow$ APAC).
  \item 100k--400k active input channels per hub.
  \item Workload: full-model inference and distributed training on 500B--1T
  parameter models.
  \item Automatic failover on fiber cut; multi-hop routing.
\end{itemize}

\textbf{New challenges addressed:}
\begin{itemize}[leftmargin=1.6em,itemsep=2pt]
  \item Multi-hop routing and wavelength reallocation.
  \item Automatic detection and failover (recovery time $<1$ second).
  \item Integration with existing cloud platforms (AWS, GCP, Azure).
  \item Regulatory compliance (GDPR, data residency laws via regional gateways).
\end{itemize}

\textbf{Cost breakdown (\$150--200M):}
\begin{itemize}[leftmargin=1.6em,itemsep=2pt]
  \item Additional optical infrastructure (2 more hubs, 2 more fiber pairs): \$80--100M
  \item Upgraded transceivers \& control systems: \$30--40M
  \item Integration with cloud platforms: \$20--30M
  \item Labor \& contingency: \$20--30M
\end{itemize}

\textbf{Success criteria:}
\begin{itemize}[leftmargin=1.6em,itemsep=2pt]
  \item Stable multi-hop routing with $<100$ ms latency for any source--destination pair
  \item Automatic failover on fiber cut: recovery time $<1$ second, no inference loss
  \item Support for 10M$+$ concurrent inferences
  \item Commercial pricing model defined (\$/Tbps-month)
  \item Open standards and reference implementation published
\end{itemize}

\textbf{Risk level: Moderate.} Requires coordination and multilateral standardization.

\subsection{Phase 3: Global Deployment (Years 7--10, \$700M--1.5B)}

\textbf{Objective:} Extend to 10+ regional hubs worldwide; become the default
infrastructure for global AI compute pooling.

\textbf{Scope:}
\begin{itemize}[leftmargin=1.6em,itemsep=2pt]
  \item Global fiber network with 10--15 datacenters across all continents.
  \item Millions of active neurons per hub.
  \item Workload: frontier models (10T$+$ parameters), multi-modal inference,
  scientific workloads.
\end{itemize}

\textbf{Cost breakdown (\$700M--1.5B):}
\begin{itemize}[leftmargin=1.6em,itemsep=2pt]
  \item New regional hubs \& fiber leases: \$400--700M
  \item Control plane, optimization, machine-learning routing: \$100--200M
  \item Integration with AI frameworks, regulatory compliance: \$150--300M
  \item Labor, testing, contingency: \$50--100M
\end{itemize}

\textbf{Projected benefits (by 2035):}
\begin{itemize}[leftmargin=1.6em,itemsep=2pt]
  \item Power savings: (200 GW baseline) $-$ (40 GW with optical) \(=\) 160 GW reduction
  (assumes 50\% adoption). Annual savings: \(\sim\!\$70\)B at typical electricity rates.
  \item Capital cost amortization: \(\sim\!\$3\)--5B/year over participating
  organizations.
  \item \textbf{Net annual benefit: \(\sim\!\$65\)B/year.}
\end{itemize}

\textbf{Success criteria:}
\begin{itemize}[leftmargin=1.6em,itemsep=2pt]
  \item Majority of global AI inference running on photonic infrastructure
  \item $<20$ ms latency for any datacenter pair
  \item Petabyte-scale daily activation transport
  \item Carbon footprint reduction of 50\% for AI compute
\end{itemize}

\textbf{Risk level: Geopolitical.} Manageable with government backing and
forward-compatible regional-fabric design.

% ======================================================================
\section{Why Cooperation Must Be Robust: Preventing Capture and Defection}
\label{sec:robust}

The infrastructure enables cooperation, but it must be designed so that
cooperation remains stable even under stress. Two scenarios matter most:

\subsection{Scenario 1: A participant attempts to defect or capture}

\textbf{The risk:} A major participant tries to seize control of the system or
unilaterally advance its own interests at others' expense.

\textbf{The mitigation:}
\begin{itemize}[leftmargin=1.6em,itemsep=2pt]
  \item \textbf{Physical distribution.} The system spans multiple jurisdictions
  and continents. No single entity controls more than a fraction of it.
  \item \textbf{Operational interdependence.} Running the system requires
  active, continuous cooperation from all major participants. If any exit or
  defect, capability degrades for all.
  \item \textbf{Verifiable routing.} Optical transport chokepoints are physical
  and auditable. Tampering with wavelength allocation or traffic patterns is
  detectable as signal loss or routing anomalies. Hidden behavior is impossible.
  \item \textbf{Governance structure.} Decisions on capacity allocation,
  maintenance, and upgrades require consensus or supermajority voting, not
  unilateral control.
\end{itemize}

\subsection{Scenario 2: A hostile actor builds a competing system}

\textbf{The risk:} An actor outside the coalition (a rogue lab, a captured
state, a misaligned actor) unilaterally defects and builds its own hostile AI,
attempting to achieve a decisive advantage.

\textbf{The mitigation:}
\begin{itemize}[leftmargin=1.6em,itemsep=2pt]
  \item \textbf{Pooled capability.} A coalitional system is, by definition, more
  capable than any single actor can build. A defector faces an opponent that is
  stronger and has more compute.
  \item \textbf{Distributed resilience.} The coalitional system survives any
  single site being compromised or attacked. A defector's system is typically
  more fragile.
  \item \textbf{Speed.} The pooled effort can respond faster to security
  threats, add defensive measures, or allocate compute to counteracting the
  hostile system.
  \item \textbf{Legitimacy and coordination.} A broadly legitimate, multi-party
  system can mobilize diplomatic, legal, and technical responses to a threat
  that a single actor cannot.
\end{itemize}

\noindent In both scenarios, the system's power comes from being \emph{cooperative}
in fact, not just by design. A participant defecting loses access to the pooled
compute and all the benefits cooperation provides. That is the enforcement
mechanism.

% ======================================================================
\section{What This Is Not: Honest Limitations}
\label{sec:limits}

\subsection{This is not a solution to alignment}

The infrastructure enables cooperation on building AI. It does nothing to make
models themselves safe, interpretable, or aligned. The hard problems of technical
AI alignment---mesa-optimization, deceptive alignment, specification gaming,
out-of-distribution robustness---remain exactly as hard as before. What the
infrastructure does is provide time and institutional structure within which that
work can proceed without race pressure.

\subsection{Efficiency is dual-use}

A 100$\times$ improvement in interconnect energy could accelerate AI scaling
rather than fund safety work. This paper proposes that the efficiency dividend
be captured for broader benefit---through governance mechanisms, benefit-sharing
commitments, and conscious choice about what problems to solve. But without that
governance, efficiency alone could make things worse.

\subsection{Shared infrastructure can become a shared chokepoint}

A unified global system with a single optical backbone could, if captured or
misgovernedy, become the ultimate centralization of power. The distributed
design, the verifiable routing, and the governance mechanisms are guards against
this. They are not guarantees. Keeping the system safe requires ongoing vigilance
and broad participation.

% ======================================================================
\section{Who This Is For: The Call to Action}
\label{sec:calls}

\subsection{For researchers and builders in AI, photonics, and systems}

You are the people who can make this real. We are asking you to:
\begin{itemize}[leftmargin=1.6em,itemsep=2pt]
  \item \textbf{Evaluate the feasibility.} Read the architecture, run the
  numbers, identify the real gaps. Where are we wrong? What have we missed?
  \item \textbf{Publish and open-source.} Design silicon-photonic reference
  implementations. Publish control-plane protocols. Build testbeds. Make this a
  public technical conversation, not a closed project.
  \item \textbf{Run Phase~1.} If the architecture is sound, the next step is
  experimental validation. Form a team. Get funding. Build it.
\end{itemize}

\subsection{For institutions that fund research (government, foundations, venture)}

You control the resources. We ask you to:
\begin{itemize}[leftmargin=1.6em,itemsep=2pt]
  \item \textbf{Treat shared AI infrastructure as strategic.} Just as the early
  internet and GPS were government-backed public goods, a planetary-scale AI
  substrate is a public good worth funding.
  \item \textbf{Co-fund Phase~1.} \$35--50M is small by the standards of AI
  development, and the returns---a demonstration that cooperation is feasible, a
  technology path that others can then scale---are enormous.
  \item \textbf{Support openness.} Fund the publication, the open-source
  implementation, and the intellectual commons that will make this a shared
  project, not a proprietary one.
\end{itemize}

\subsection{For AI labs and companies}

You have the most to gain and the most to lose. We ask you to:
\begin{itemize}[leftmargin=1.6em,itemsep=2pt]
  \item \textbf{Do the math.} Compute the cost of racing. Compute the cost of
  pooling. Compute the capability you could build cooperatively. Then ask
  yourself: which is better for your business, your people, and your role in
  the world?
  \item \textbf{Form a consortium.} Announce a commitment to cooperation on
  shared infrastructure. The first mover gets outsized credit; the second mover
  looks like they were forced into it.
  \item \textbf{Open the technical work.} You do not need to give away your
  models or your data. But you can contribute to the photonic fabric, the
  protocols, and the standards that everyone will eventually use anyway. Get
  ahead of the curve.
\end{itemize}

\subsection{For policymakers and diplomats}

You set the rules of the
\begin{itemize}[leftmargin=1.6em,itemsep=2pt] \item \textbf{Recognize the opportunity.} A cooperatively built global AI system is not a loss of sovereignty; it is a multiplication of it. Each participant retains full control of its own compute and data. What is shared is transport and infrastructure---the pipes, not the content. \item \textbf{Enable cooperation.} Negotiate fiber-access agreements. Work with cable consortiums to unlock dark fiber for AI. Remove regulatory barriers to shared infrastructure. \item \textbf{Govern transparently.} The governance of the system itself should be transparent, multi-stakeholder, and legitimately representative. Use this as a model for how to govern other powerful technologies. \end{itemize}
\subsection{For everyone else}

You care about the future. You want AI to solve hard problems rather than
exacerbate them. You want a world where capability is broadly distributed rather
than hoarded. \textbf{Talk about this.} Demand that the actors above treat it
seriously. Make it clear that cooperation is what you want to see, that you will
support it, and that you will hold builders accountable if they choose to race
instead.

% ======================================================================
\section{The Broader Point}
\label{sec:broader}

This proposal is about more than infrastructure. It is about \emph{choice}.

Right now, the development of powerful AI is being decided by default---through
competition, through the logic of markets and geopolitical rivalry, through the
dynamics of venture capital and research funding. That may or may not produce
good outcomes. We will find out only after the fact.

This proposal asks: what if we chose instead? What if, while we still can, we
decided together what we want AI to be---not hostile, not captured, not
optimizing for the private benefit of a few actors, but genuinely useful,
broadly distributed, and built for the problems that matter?

The photonic network is not the point. The point is that it makes choice
possible. With the data-movement bottleneck removed, cooperation becomes not a
nice idea but a practical engineering path. One that is faster, more capable,
less wasteful, and more useful than the alternatives.

We are asking you to make that choice deliberately. To pool resources, to work
openly, to govern fairly, and to build something together that no one could
build alone. Not because it is virtuous---though it is---but because it is the
best engineering, the fastest path to capability, and the most likely to
actually help.

The window for this choice is closing. In a few years, the systems being built
will be powerful enough that they become hard to redirect. The institutions will
be locked in. The geopolitical positions will be calcified. But right now---today,
this year, in the next few years---it is still possible to choose.

We invite you to choose together.

% ======================================================================
\section*{Acknowledgments}

This manuscript benefited from substantial drafting, synthesis, architectural
framing, and technical modeling assistance from a large language model (Anthropic
Claude 4.5 Haiku). The human author is responsible for all claims, framing,
strategic positioning, and errors. Feedback, criticism, and collaboration are
explicitly welcomed.

The empirical and philosophical foundations for the cooperative approach were
developed in parallel work, including the Empirical Pascal's Wager framework and
related AI governance analysis by the author.

% ======================================================================
\section*{Author Contributions}

[Author Name] conceived the proposal, framed the arguments for a cooperative
approach, and oversaw all drafting and revision. Technical architecture and
performance modeling were developed collaboratively with AI assistance. This is
a working pre-print intended for circulation, feedback, and critique from the
Georgia Tech AI Safety Initiative, the photonics and distributed-systems
research community, and potential consortium partners.

% ======================================================================
\section*{License and Availability}

This work is released under the Creative Commons Attribution 4.0 International
License (CC-BY-4.0). Source, revision history, and collaboration materials are
maintained on GitHub at [repository URL, to be filled in].

% ======================================================================
\begin{thebibliography}{99}

\bibitem{armstrong2016racing}
S.~Armstrong, N.~Bostrom, and C.~Shulman,
\newblock ``Racing to the Precipice: A Model of Artificial Intelligence Development,''
\newblock \emph{AI \& Society}, vol.~31, no.~2, pp.~201--206, 2016.

\bibitem{dafoe2018agenda}
A.~Dafoe,
\newblock ``AI Governance: A Research Agenda,''
\newblock Centre for the Governance of AI, Future of Humanity Institute, University of Oxford, 2018.

\bibitem{openai2018charter}
OpenAI,
\newblock ``OpenAI Charter,''
\newblock \url{https://openai.com/charter/}, 2018.

\bibitem{bostrom2014superintelligence}
N.~Bostrom,
\newblock \emph{Superintelligence: Paths, Dangers, Strategies},
\newblock Oxford University Press, 2014.

\bibitem{okeefe2020windfall}
C.~O'Keefe, T.~Anderljung, and S.~Legg,
\newblock ``Should Artificial Intelligence Governance be Centralised? Design Lessons from History,''
\newblock In \emph{Proceedings of the AAAI/ACM Conference on AI, Ethics, and Society}, 2020.

\bibitem{amodei2024grace}
D.~Amodei and D.~Amodei,
\newblock ``Machines of Loving Grace,''
\newblock OpenAI Blog, 2024.

\bibitem{shen2017deep}
Y.~Shen, N.~C.~Harris, S.~Skirlo, M.~Prabhu, T.~Baehr-Jones, M.~Hochberg, X.~Sun, S.~Zhao, H.~Larochelle, D.~Englund, and M.~Soljavc,
\newblock ``Deep Learning with Coherent Nanophotonic Circuits,''
\newblock \emph{Nature Photonics}, vol.~11, pp.~441--446, 2017.

\bibitem{wetzstein2020inference}
G.~Wetzstein, A.~Ozcan, S.~Gigan, S.~Fan, D.~Englund, M.~Soljavc, C.~Denz, D.~A.~B.~Miller, and D.~Psaltis,
\newblock ``Inference in Artificial Intelligence with Deep Optics and Photonics,''
\newblock \emph{Nature}, vol.~588, pp.~39--47, 2020.

\bibitem{feldmann2021parallel}
J.~Feldmann, N.~Youngblood, M.~Karpov, et al.,
\newblock ``Parallel Convolutional Processing Using an Integrated Photonic Tensor Core,''
\newblock \emph{Nature}, vol.~589, pp.~52--58, 2021.

\bibitem{patterson2021carbon}
D.~Patterson, J.~Gonzalez, Q.~Le, C.~Liang, L.-M.~Munguia, D.~Rothchild, D.~So, M.~Texier, and J.~Dean,
\newblock ``Carbon Emissions and Large Neural Network Training,''
\newblock \emph{arXiv:2104.10350}, 2021.

\bibitem{devries2023energy}
A.~de Vries,
\newblock ``The Growing Energy Footprint of Artificial Intelligence,''
\newblock \emph{Joule}, vol.~7, no.~10, pp.~2191--2194, 2023.

\bibitem{iea2024electricity}
International Energy Agency,
\newblock ``Electricity 2024: Analysis and Forecast to 2026,''
\newblock IEA, Paris, 2024.

\bibitem{sastry2024compute}
G.~Sastry, L.~Heim, H.~Belfield, M.~Anderljung, M.~Brundage, et al.,
\newblock ``Computing Power and the Governance of Artificial Intelligence,''
\newblock \emph{arXiv:2402.08797}, 2024.

\bibitem{brundage2020verifiable}
M.~Brundage, S.~Avin, J.~Wang, H.~Belfield, et al.,
\newblock ``Toward Trustworthy AI Development: Mechanisms for Supporting Verifiable Claims,''
\newblock \emph{arXiv:2004.07213}, 2020.

\bibitem{hausenloy2023magic}
R.~Hausenloy, E.~Anderljung, and A.~Andersson,
\newblock ``Multinational AGI Consortium (MAGIC),''
\newblock \emph{Differential AI Governance}, 2023.

\bibitem{achesonlilienthal1946}
D.~E.~Acheson and B.~Lilienthal,
\newblock ``A Report on the International Control of Atomic Energy,''
\newblock U.S. Department of State, 1946.

\bibitem{winzer2018optical}
P.~J.~Winzer and D.~T.~Neilson,
\newblock ``From Scaling Chaos to the Roadmap of Lightwave Communication,''
\newblock \emph{Nature Photonics}, vol.~12, no.~6, pp.~303--314, 2018.

\bibitem{telegeography2023submarine}
TeleGeography,
\newblock ``TeleGeography Submarine Cable Map,''
\newblock \url{https://www2.telegeography.com/submarine-cable-faqs-frequently-asked-questions}, 2023.

\end{thebibliography}

\end{document}
